\documentclass[12pt, a4paper]{ctexart}
\usepackage{fancyhdr, amsmath, amssymb, geometry, enumitem}
\geometry{top=2.5cm, bottom=2.5cm, left=2cm, right=2cm}
\setCJKmainfont{Hiragino Sans GB}
\setCJKsansfont{Heiti SC}
\pagestyle{fancy}
\fancyhf{}
\fancyhead[L]{\small 培优专题·代数变形技巧}
\fancyhead[R]{\small 测试卷}
\fancyfoot[C]{\thepage}
\renewcommand{\headrulewidth}{0.4pt}

\begin{document}

\begin{center}
{\Large\bfseries 培优专题·代数变形技巧\quad 测试卷}\\[6pt]
{\normalsize 总分：150分\quad 考试时间：120分钟}
\end{center}
\vspace{6pt}
\noindent\rule{\textwidth}{0.5pt}

\section*{一、单项选择题（本大题共8小题，每小题5分，共40分。每小题只有一个正确选项。）}

\begin{enumerate}

\item 若$x+\dfrac{1}{x}=3$，则$x^2+\dfrac{1}{x^2}$的值为\ （\quad）
\begin{enumerate}[label=\Alph*., itemsep=0pt, topsep=2pt]
\item $7$ \quad \item $9$ \quad \item $11$ \quad \item $5$
\end{enumerate}

\item 已知$a+b+c=0$，$a^2+b^2+c^2=1$，则$ab+bc+ca$的值为\ （\quad）
\begin{enumerate}[label=\Alph*., itemsep=0pt, topsep=2pt]
\item $1$ \quad \item $-1$ \quad \item $-\dfrac{1}{2}$ \quad \item $\dfrac{1}{2}$
\end{enumerate}

\item 将$x^4+4$因式分解（在有理数范围内），结果为\ （\quad）
\begin{enumerate}[label=\Alph*., itemsep=0pt, topsep=2pt]
\item $(x^2+2x+2)(x^2-2x+2)$ \quad \item $(x^2+2)(x^2-2)$
\item $(x^2+2x-2)(x^2-2x-2)$ \quad \item $(x^2+2x+2)(x^2+2x-2)$
\end{enumerate}

\item 已知$x$、$y$、$z$满足$x+y+z=3$，$xy+yz+zx=3$，则$x^3+y^3+z^3-3xyz$的值为\ （\quad）
\begin{enumerate}[label=\Alph*., itemsep=0pt, topsep=2pt]
\item $0$ \quad \item $3$ \quad \item $9$ \quad \item $27$
\end{enumerate}

\item 设$a$、$b$为实数，且$a^3+b^3=2$，则$a+b$的最大值为\ （\quad）
\begin{enumerate}[label=\Alph*., itemsep=0pt, topsep=2pt]
\item $1$ \quad \item $2$ \quad \item $\sqrt[3]{4}$ \quad \item $\sqrt{2}$
\end{enumerate}

\item 已知多项式$f(x)=x^3-3x+1$的三个实根为$x_1$、$x_2$、$x_3$，则$\dfrac{1}{x_1}+\dfrac{1}{x_2}+\dfrac{1}{x_3}$的值为\ （\quad）
\begin{enumerate}[label=\Alph*., itemsep=0pt, topsep=2pt]
\item $0$ \quad \item $3$ \quad \item $-3$ \quad \item $1$
\end{enumerate}

\item 若$a+b=1$，$a^2+b^2=2$，则$a^5+b^5$的值为\ （\quad）
\begin{enumerate}[label=\Alph*., itemsep=0pt, topsep=2pt]
\item $\dfrac{11}{4}$ \quad \item $\dfrac{19}{4}$ \quad \item $\dfrac{7}{2}$ \quad \item $\dfrac{9}{2}$
\end{enumerate}

\item 将$x^6-1$在实数范围内因式分解，不同的一次因式（含常数因子）共有\ （\quad）
\begin{enumerate}[label=\Alph*., itemsep=0pt, topsep=2pt]
\item $4$个 \quad \item $3$个 \quad \item $2$个 \quad \item $6$个
\end{enumerate}

\end{enumerate}

\section*{二、多项选择题（本大题共4小题，每小题5分，共20分。每小题有多个正确选项，全部选对得5分，部分选对得2分，有选错得0分。）}

\begin{enumerate}[resume]

\item 已知$a+b+c=0$，下列结论正确的是\ （\quad）
\begin{enumerate}[label=\Alph*., itemsep=0pt, topsep=2pt]
\item $a^3+b^3+c^3=3abc$
\item $a^2+b^2+c^2=-2(ab+bc+ca)$
\item $\dfrac{1}{a^2}+\dfrac{1}{b^2}+\dfrac{1}{c^2}=\left(\dfrac{1}{a}+\dfrac{1}{b}+\dfrac{1}{c}\right)^2$
\item $a^5+b^5+c^5=5abc(ab+bc+ca)$
\end{enumerate}

\item 关于因式分解，下列说法正确的是\ （\quad）
\begin{enumerate}[label=\Alph*., itemsep=0pt, topsep=2pt]
\item $x^4+y^4=(x^2+y^2+\sqrt{2}xy)(x^2+y^2-\sqrt{2}xy)$
\item $a^3+b^3+c^3-3abc=(a+b+c)(a^2+b^2+c^2-ab-bc-ca)$
\item $x^n-y^n$对任意正整数$n$都可被$x-y$整除
\item $x^n+y^n$对任意正奇数$n$都可被$x+y$整除
\end{enumerate}

\item 设$p$、$q$为实数，多项式$x^2+px+q$的两个根为$x_1$、$x_2$，则下列等式正确的是\ （\quad）
\begin{enumerate}[label=\Alph*., itemsep=0pt, topsep=2pt]
\item $x_1+x_2=-p$
\item $x_1^2+x_2^2=p^2-2q$
\item $x_1^3+x_2^3=-p^3+3pq$
\item $|x_1-x_2|=\sqrt{p^2-4q}$
\end{enumerate}

\item 已知$\omega$为$1$的三次虚根（即$\omega^3=1$，$\omega\neq 1$），下列结论正确的是\ （\quad）
\begin{enumerate}[label=\Alph*., itemsep=0pt, topsep=2pt]
\item $\omega^2+\omega+1=0$
\item $\omega^2=\bar{\omega}$
\item $(1-\omega)(1-\omega^2)=3$
\item $1+\omega+\omega^2+\omega^3+\omega^4+\omega^5=0$
\end{enumerate}

\end{enumerate}

\section*{三、填空题（本大题共4小题，每小题5分，共20分。）}

\begin{enumerate}[resume]

\item 计算$(1+x)(1+x^2)(1+x^4)(1+x^8)$，用$x$的幂次表示为$\dfrac{1-x^n}{1-x}$，则$n=$\underline{\qquad}。

\item 已知$a+b=1$，$a^2+b^2=2$，则$a^4+b^4=$\underline{\qquad}。

\item 设$x_1$、$x_2$是方程$x^2-3x+1=0$的两个根，则$x_1^4+x_2^4=$\underline{\qquad}。

\item $x^5-1$在有理数范围内可分解为$(x-1)(x^4+x^3+x^2+x+1)$。$x^4+x^3+x^2+x+1$在实数范围内可分解为两个二次因式的乘积：$(x^2+ax+1)(x^2+bx+1)$，其中$a=$\underline{\qquad}，$b=$\underline{\qquad}。

\end{enumerate}

\section*{四、解答题（本大题共6小题，共70分。解答应写出文字说明、证明过程或演算步骤。）}

\begin{enumerate}[resume]

\item （本小题满分10分）\\
已知数列$\{a_n\}$满足$a_1=1$，$a_{n+1}=2a_n+3$（$n\geq 1$）。
\begin{enumerate}[(1)]
\item 设$b_n=a_n+3$，证明$\{b_n\}$为等比数列，并求$a_n$的通项公式；
\item 求$a_1+a_2+\cdots+a_n$。
\end{enumerate}

\item （本小题满分10分）\\
已知$x+y+z=1$，$x^2+y^2+z^2=2$，$x^3+y^3+z^3=3$。
\begin{enumerate}[(1)]
\item 求$xyz$的值；
\item 求$x^4+y^4+z^4$的值。
\end{enumerate}

\item （本小题满分12分）\\
已知正整数$n\geq 2$。
\begin{enumerate}[(1)]
\item 证明$x^n-1=(x-1)(x^{n-1}+x^{n-2}+\cdots+x+1)$；
\item 当$n$为奇数时，证明$x^n+1$可被$x+1$整除；
\item 将$x^6-1$在实数范围内完全因式分解。
\end{enumerate}

\item （本小题满分12分）\\
已知$S=\cos\dfrac{\pi}{7}\cos\dfrac{2\pi}{7}\cos\dfrac{3\pi}{7}$。
\begin{enumerate}[(1)]
\item 利用$\sin\dfrac{\pi}{7}$将$S$化简，并证明$S=\dfrac{1}{8}$；
\item 利用$x^7-1=0$的根与韦达定理，验证上述结论。
\end{enumerate}

\item （本小题满分13分）\\
已知$a$、$b$、$c$为正实数，$abc=1$。
\begin{enumerate}[(1)]
\item 证明$\dfrac{1}{a}+\dfrac{1}{b}+\dfrac{1}{c}\geq 3$；
\item 证明$a^3+b^3+c^3\geq a+b+c$；
\item 证明$(a+b)(b+c)(c+a)\geq 8$。
\end{enumerate}

\item （本小题满分13分）\\
已知多项式$P(x)=x^3+ax^2+bx+c$的三个根为$x_1$、$x_2$、$x_3$，且$x_1^2+x_2^2+x_3^2=10$，$x_1^3+x_2^3+x_3^3=30$。
\begin{enumerate}[(1)]
\item 若$x_1+x_2+x_3=1$，求$a$、$b$、$c$的值；
\item 求$x_1^4+x_2^4+x_3^4$的值；
\item 证明：$x_1$、$x_2$、$x_3$中至少有一个大于$2$。
\end{enumerate}

\end{enumerate}

\end{document}
